% !TEX root = intro.tex
\paragraph{Human-AI collaboration.} 
The nature of human-AI collaboration can be broadly classified into two types: \emph{augmentation} and \emph{automation} \citep{brynjolfsson2017can, raisch2021artificial}. In particular, augmentation represents a ``human-in-the-loop'' type workforce where human experts combine machine learning with their own judgement to make better decisions. Such augmentation is often found in high-stake settings such as healthcare \citep{lebovitz2022engage}, child maltreatment hotline screening \citep{De-ArteagaFC20}, refugee resettlement \citep{bansak2020outcome,ahani2023dynamic} and bail decisions \citep{kleinberg2018human}. A general concern for augmentation relates to humans' compliance patterns \citep{bastani2021improving,lebovitz2022engage}, the impact of such patterns to decision accuracy \citep{kleinberg2018human,mclaughlin2022algorithmic} and the impact on fairness \citep{McLaughlinSG22,ge2023rethinking}. Automation, on the other hand, concerns the use of machine learning in place of humans and is widely applied in human and social services where colossal demands overwhelm limited human capacity, such as in data labelling \citep{vishwakarma2023promises}, content moderation \citep{gorwa2020algorithmic,makhijani2021quest} and insurance \citep{eubanks2018automating}. Full automation is clearly undesirable in these applications as machine learning can err. A natural question is how one can better utilize the limited human capacity when machine is uncertain about the prediction. The literature on ``learning to defer'', which studies when machine learning algorithms should defer decisions to downstream experts, tries to answer this question and is where our work fits in. 

Assuming that humans have perfect prediction ability but there is limited capacity, classical learning to defer has two streams of research, \emph{learning with abstention} and \emph{selective sampling}. Although these two streams have been studied separately (with a few exceptions, see discussion below), our work provides an endogenous approach to connect them.  In particular, learning with abstention can be traced back to \cite{Chow57,Chow70} and studies an offline classification problem with the option to not classify a data point for a fixed cost. For example, in content moderation, this corresponds to paying a fixed fee to an exogenous human reviewer to review a post. Since optimizing the original problem is computational infeasible, an extensive line of work investigates suitable surrogate loss functions and optimization methods; see \cite{BartlettW08,El-YanivW10,CortesDM16a} and references therein. For online learning with expert advice, it is shown that even a limited amount of abstention allows better regret bound than without \citep{SayediZB10,LiLWS11,ZhangC16,NeuZ20}; \cite{CortesDGMY18} studies a similar problem but allows experts to also abstain. Selective sampling (or label efficient prediction) considers a different model where the algorithm makes predictions for every arrival; but the ground truth label is unavailable unless the algorithm queries for it; if queried, the label is available immediately \citep{Cesa-BianchiLS05}. The goal is to obtain low regret while using as few queries as possible. Typical solutions query only when a confidence interval exceeds a certain threshold \citep{Cesa-BianchiLS05,cesa2009robust,OrabonaC11,DekelGS12}. There is recent work connecting the two directions by showing that allowing abstention of a small fixed cost can lead to better regret bound for selective sampling \citep{Zhu022a,PuchkinZ22}. However, these works treat the abstention cost and query limits exogenously, neglecting  the impact of limited human capacity to learning to defer \citep{Leitao22}. Our model serve as one step to capture it by endogenously connecting both costly abstention and selective sampling via delays. In particular, our admission component determines both abstention and selective sampling. In addition, the cost of abstention is dynamically affected by the delay in getting human reviews whereas the avoidance of frequent sampling (to get data) is captured by its impact on the delay. 

Although the above line of work as well as ours assumes perfect labels from human predictions, a more recent stream of work on ``learning to defer'' considers imperfect human predictions and is focused on combining prediction ability of experts and learning algorithms; this moves towards the augmentation type of human-machine collaboration. In particular, \cite{MadrasCPZ18, MozannarS20,wilder2020learning,CharusaieMSS22} study settings with an offline dataset and expert labels. The goal is to learn both a classifier, which predicts outcome, and a rejector that predicts when to defer to human experts. The loss is defined by the machine's classification loss over non-deferred data, humans' classification loss over deferred data, and the cost to query experts. \cite{DeKGG20,DeOZR21} study a setting where given expert loss functions, the algorithm picks a size-limited subset of data to outsource to humans and solves a regression or a classification problem on remaining data, with a goal to minimize the total loss. \cite{raghu2019algorithmic} extends the model by allowing the expert classification loss to depend on human effort and further considering an allocation of human effort to different data points. \cite{KeswaniLK21,VermaBN23,mao2023principled} consider learning to defer with multiple experts. 

\paragraph{Bandits with knapsacks or delays.} Restricting our attention to admission decisions, our model bears similar challenges with the literature on bandits with knapsacks or broadly online learning with resource constraints, which finds applications in revenue management \citep{besbes2012blind,wang2014close,ferreira2018online}. In particular, for bandits with knapsacks, there are arrivals with rewards and required resources from an unknown distribution. The algorithm only observes the reward and required resources after admitting an arrival and the goal is to obtain as much reward as possible subject to resource constraints \citep{badanidiyuru2018bandits,agrawal2019bandits}. A typical primal-dual approach learns the optimal dual variables of a fixed fluid model and explores with upper confidence bound \citep{badanidiyuru2014resourceful,LiSY21}; these extend to contextual settings \citep{WuSLJ15,AgrawalD16,AgrawalDL16, SlivkinsSF23}, general linear constraints \citep{PacchianoGBJ21, LiuLSY21} and constrained reinforcement learning \citep{BrantleyDLMSSS20,Cheung19}. These results cannot immediately apply to our setting for two reasons. First, the resource constraint in our model is dynamically captured by time-varying queue lengths instead of a single resource constraint over the entire horizon; thus learning fixed dual variables is insufficient for good performance. Second, our admission decisions must also consider the effect on classification and relying only on optimism-based exploration is inefficient (see Section~\ref{sec:opti-fail}).

The problem of bandits with delays is related to our setting where feedback of an admitted job gets delayed due to congestion. Motivated by conversion in online advertising, bandits with delays consider the problem where the reward of each pulled arm is only revealed after a random delay independently generated from a fixed distribution \citep{DudikHKKLRZ11,Chapelle14}. Assuming independence between rewards and delays as well as bounded delay expectation, \cite{JoulaniGS13,MandelLBP15} propose a general reduction from non-delay settings to their delayed counterpart. Subsequent papers consider censored settings with unobservable rewards \citep{VernadeCP17}, general (heavy-tail) delay distributions \citep{ManegueuVCV20, WuW22} and reward-dependent delay \citep{LancewickiSKM21}. \cite{VernadeCLZEB20,blanchet2023delay} study (generalized) linear contextual bandit with delayed feedback. The key difference between bandits with delays and our model is that delays for jobs in our model are not independent across jobs due to queueing effect; we provide a more elaborate comparison to those works in Appendix~\ref{app:comparison}. 

\paragraph{Learning in queueing systems.} Learning in queueing systems can be classified into two types: 1) learning to schedule with unknown service rates to obtain low delay; 2) learning unknown utility of jobs / servers to obtain high reward in a congested system. For the first line of research, an intuitive approach to measure delay suboptimality is via the queueing regret, defined as the difference in queue lengths compared with a near-optimal algorithm \citep{walton2021learning}. The interest for queueing regret is in its asymptotic scaling in the time horizon, and it is studied for single-queue multi-server systems \citep{Walton14,KrishnasamySJS21,StahlbuhkSM21}, multi-queue single-server systems \citep{krishnasamy2018learning}, load balancing \citep{choudhury2021job}, queues with abandonment \citep{zhong2022learning} and more general markov decision processes with countable infinite state space \citep{adler2023bayesian}. As an asymptotic metric may not capture the learning efficiency of the system, \cite{freund2023quantifying} considers an alternative metric (cost of learning in queueing) that measures transient performance by the maximum increase in time-averaged queue length. This metric is motivated by works that study stabilization of queueing systems without knowledge of parameters. In particular, \cite{NeelyRP12,yang2023learning,nguyen2023learning} combine the celebrated MaxWeight scheduling algorithm \citep{tassiulas1992stability} with either discounted UCB or sliding-window UCB for scheduling with time-varying service rates. \cite{FuHL22,gaitonde2023price} study decentralized learning with strategic queues and \cite{StahlbuhkSM19,sentenac2021decentralized,FreundLW22} consider efficient decentralized learning algorithms for cooperative queues. Although most work for learning in queues focus on an online stochastic setting, \cite{huang2023queue,liang2018minimizing} study online adversarial setting and \cite{singh2022feature} considers an offline feature-based setting.

Our work is closer to the second literature that learns job utility in a queueing system. In particular, \cite{massoulie2018capacity, ShahGMV20} consider Bayesian learning in an expert system where jobs are routed to different experts for labels and the goal is to keep the expert system stable. \cite{johari2021matching,hsu2022integrated,fu2022optimal} study a matching system where incoming jobs have uncertain payoffs when served by different servers and the objective is to maximize the total utility of served jobs within a finite horizon. \cite{jia2022online,Jia0S22,chen2023online} investigate regret-optimal learning algorithms and \cite{li2023experimenting} studies randomized experimentation for online pricing in a queueing system. 

We note that, although most performance guarantees for learning in queueing systems deteriorate as the number of job types $K$ increases, our work allows admission, scheduling and learning in a many-type setting where the performance guarantee is independent of $K$. Although prior work obtains such a guarantee in a Bayesian setting, where the type of a job corresponds to a distribution over a finite set of labels \citep{argon2009priority,massoulie2018capacity,ShahGMV20}, their service rates are only server-dependent (thus finite). In contrast, our work allows for job-dependent service rates. In addition, a Bayesian setting does not immediately capture the contextual information between jobs that may be useful for learning. We note that \cite{singh2022feature} consider a multi-class queueing system where a job has an observed feature vector and an unobserved job type, and the task is to assign jobs to a fixed number of classes with the goal of minimizing mean holding cost, with known holding cost rate for any type. They find that directly optimizing a mapping from features to classes can greatly reduce the holding cost, compared with a predict-then-optimize approach. Different from their setting, we consider an online learning setting where reviewing a job type provides information for other types. This creates an explore-exploit trade-off complicated with the additional challenge that the feedback experiences queueing delay. To the best of our knowledge, our work is the first result for efficient online learning in a queueing system with contextual information.
 
\paragraph{Joint admission and scheduling.} When there is no learning, our problem becomes a joint admission and scheduling problem that is widely studied in wireless networks and the general focus is on a decentralized system \citep{kelly1998rate,lin2006tutorial}. Our method is based on the drift-plus penalty algorithm \citep{neely2022stochastic}, which is a common approach for joint admission and scheduling, first noted as a greedy primal-dual algorithm in \cite{stolyar2005maximizing}. The intuition is to view queue lengths as dual variables to guide admission; \cite{huang2009delay} formalizes this idea and exploits it to obtain better utility-delay tradeoffs. 